
\begin{abstract_online}{Assessment Tools in Maple: Recent Developments and Challenges}{%
        \underline{Paulina Chin}$^{1}$, Paul DeMarco$^{1}$}{[pchin@maplesoft.com]
        }{%
        $^1$ Maplesoft, Waterloo, Ontario, Canada}

The benefits of using computer algebra systems to  demonstrate and explore mathematical concepts are clear. However, the use of a CAS for assessment in mathematics and science courses still poses a number of challenges.  In this presentation, we will show several recent additions to the Maple [1] software package for the purposes of grading and self-assessment, and our focus will be on the design issues we encountered in building the tools as well as the challenges we face in their future development.  

One of these tools is the graph assessment tool in Maple’s Grading package. Its purpose is to allow “sketches” of plots entered by students into a computer to be compared to ones requested by an instructor. Issues that we needed to consider in our design that continue to pose difficulties include noise in the data points, scaling of test questions, and questions that do not have unique solutions. 

We will also present and discuss the Quiz command in the Grading package, which is intended to generate randomized quizzes on a variety of mathematical subjects, with automated assessment of the answers. The challenges in the design of this tool include going beyond simple multiple-choice questions and providing an easy-to-use interface that allows instructors to focus on the mathematical concepts rather than programming-like syntax for authoring the questions.

Finally, we will show the EssayTools package, which contains tools related to the assessment of essays. Though they are not meant for the assessment of mathematics, the algorithms were easily implemented using the functionality a CAS offers.
     

\textbf{Keywords} \newline{} software, assessment, education, Maple

\textbf{References} \newline{}
% Model for references to books.
[1] Maple 2019, www.maplesoft.com/products/maple \newline{}
\end{abstract_online}
    